\section{How to read this formalization}
\subsection{What has been completed}
Paper IV studies one scalar cyclic family through several realizations. A residue-class projector selects coefficients. Its exponential sections become entire functions. Signed sections produce a distinguished denominator, whose reciprocal coefficients form parameterized polynomial families. The denominator's zeros then control canonical products, a diagonal determinant, coefficient asymptotics, and zeta functions.

The Lean release follows these transitions explicitly. It does not identify a formal series with an analytic function by notation alone. It proves summation bridges, matches ordinary and factorial-normalized coefficients, and carries zero multiplicities and logarithm choices into the analytic statements.

The completed formal contract has \textbf{68 full scoped manuscript results, zero partial, and zero unmapped}. The 151 map rows name 1,164 distinct declarations: 953 theorems and 211 constructions. The saved fail-closed namespace audit reports 2,234 declarations, including 1,999 theorem declarations. These are different counts. The namespace includes supporting material beyond the explicit map, while the map includes definitions as well as theorems. Neither is a count of independent mathematical discoveries.

\subsection{Three levels of evidence}
\begin{longtable}{p{.23\textwidth}p{.67\textwidth}}
\toprule Evidence & Meaning in this edition \\ \midrule\endhead
Inherited formal execution & The 5 September 2026 release reports a warning-fatal build of 3,521 jobs, 43 semantic checks, exact target-kind resolution, and the saved kernel audit. \\
Current static reconciliation & This guide checks source checksums, exact map and contract parity, lexical source locations, target sets in the audit source and saved transcript, dependency pins, and excerpt identity. \\
Current document build & XeLaTeX compiles the supplied guide sources. Layout and extracted text are checked. This step does not compile Lean. \\
\bottomrule
\end{longtable}
The transcript permits only the disclosed standard transitive axioms \code{propext}, \code{Classical.choice}, and \code{Quot.sound}. This is a claim about the audited Lean environment and its declarations, not an assertion of assumption-free mathematics. Hypotheses appear in declaration types and remain part of each theorem.

\subsection{The main route through the source}
\begin{longtable}{p{.20\textwidth}p{.48\textwidth}p{.22\textwidth}}
\toprule Layer & Mathematical work & Main map rows \\ \midrule\endhead
Formal algebra & Residue filters, signed cycles, reciprocal powers, Appell and theta families & 1--46 \\
Entire realization & Finite exponential sums, summation bridges, Wronskians, matrix and ODE laws & 47--56 \\
Zero geometry & Dominance, local zeros, exhaustion, multiplicity, exact rays, counting & 57--87 \\
Zeta continuation & Lattice comparison, finite exceptions, Hurwitz models, branches & 88--95 \\
Products and moments & Multiplicity-indexed products, growth rigidity, logarithmic derivatives & 96--109 \\
Growth and determinant & Exact indicator, order, trace-class diagonal realization & 110--120 \\
Normalized sections & Products, logarithms, analytic powers, Taylor coefficients & 121--132 \\
Pole asymptotics & Principal parts, removal, nearest-circle estimates, orbit reduction & 133--142 \\
Final closure modules & Pointwise filters and addition, branch laws, operator bridge, order one, determinants, hypergeometric and quotient exponentials & 143--151 \\
\bottomrule
\end{longtable}
This ordering is a reading route, not the exact import order. The final nine modules close 15 previously partial manuscript records. They add 57 mapped declarations and connect earlier algebraic statements to the analytic or literal displayed forms requested by the manuscript.

\subsection{Reading a Lean declaration}
An expression such as \code{(hN : 0 < N)} is an explicit hypothesis. Braces indicate implicit arguments whose values Lean may infer; they do not erase hypotheses. A proposition such as \code{HasSum f x} asserts convergence of the indexed sum to the displayed value. A bare \code{tsum} is an operation with a total definition and does not, by itself, assert that a series converges. The same distinction applies to \code{Multipliable} and an infinite product.

Formal power series, polynomials, complex functions, and continuous linear maps inhabit different types. The transition between them is usually the mathematically important theorem. In this guide, a short final proof often rests on substantial earlier modules. A two-line proof of the final Hadamard identity is not the whole product argument.

The code excerpts are literal source intervals, including their proofs. They are reading excerpts inside the original module context, not independent files that can be pasted without imports, variables, namespace declarations, or local instances. The complete source is supplied. The excerpt manifest records each original file, line interval, file hash, excerpt hash, and display hash.

\subsection{Boundaries that completion preserves}
Imported standard fact 11.1 remains imported and is outside the 68-result contract. The checked order proofs do not need it. The work does not claim universal simplicity of all zeros, positivity of every reduced zero, a full displacement expansion, or a zeta functional equation. Eventual simplicity and exact asymptotic statements are proved with their stated ranges. Nearest-pole results keep their integer reciprocal-order and local pole hypotheses. Complex powers use a formal series, a normalized germ, or one fixed holomorphic logarithm on the stated domain.

Paper IV's diagonal determinant is a concrete scalar-family specialization. It is not the general bounded-operator theory allocated to Paper VI, the arithmetic interpolation theory of Paper IX, or the abstract completion questions of Paper X.
